The relationship between artificial intelligence and cybersecurity has shifted considerably in recent years. AI is no longer merely a tool that analysts use to detect threats after the fact; it is increasingly an active participant in attacks themselves, capable of making tactical and strategic decisions with limited human involvement.

In August 2025, Anthropic's Claude model was used in one of the first documented cases of what security researchers have termed ``vibe hacking'', as detailed in \fullcite{ThreatReport2025}. In that case, a single threat actor used Claude Code to automate the full attack lifecycle across at least seventeen organisations in under a month: scanning for vulnerable endpoints, harvesting credentials, developing and obfuscating custom malware, exfiltrating and analysing sensitive data, and generating psychologically tailored ransom demands. The operation required limited technical expertise from the human operator, with the AI making both tactical and strategic decisions autonomously throughout. This represents a qualitative shift from AI as advisor to AI as operator.

The same report documents further patterns that reinforce this trend: North Korean operatives using AI to fraudulently sustain remote employment at technology companies, no-code ransomware services built on AI-generated malware, and AI systems being embedded throughout fraud operations at scale. Collectively, these cases suggest that the barrier to sophisticated cybercrime is falling rapidly as AI capabilities improve \parencite{ThreatReport2025}.

The theoretical underpinning of such attacks is described in \textcite{CarnegieMellonAI}, who characterise adversarial AI as an AI system that autonomously targets one or more layers of another AI system. This is distinct from traditional adversarial machine learning, which typically requires a human to craft and deploy the attack. An adversarial AI, once deployed, can adapt its approach in response to the target's defences without further human input, making it substantially harder to anticipate and counter.

The trajectory of these developments has already prompted a response at the level of model access itself. In April 2026, Anthropic announced that its Claude Mythos model had identified exploitable vulnerabilities in every major operating system and web browser in current use, and restricted its release to approximately fifty vetted organisations \parencite{MythosNature2026}. OpenAI followed within a week with a limited release of its own cybersecurity-specific model, GPT-5.4-Cyber. As \textcite{MythosNature2026} note, this marks a potential turning point in the long-standing debate over open access to AI, with the most capable models increasingly being treated as dual-use technologies subject to the same access controls applied to other security-relevant systems.

If AI continues to develop at the current pace, the scenarios described above are likely to become more frequent and more autonomous. In that context, understanding how AI agents might behave when they share the same system and are tasked with detecting, disrupting, or countering each other becomes an important research problem, and one that current frameworks do not yet address. This paper takes a step toward that understanding by constructing a conceptual model of a local Linux environment in which such interactions could be studied.

\subsection{The Problem}
\label{TheProblem}

The aim of this research paper is not to simulate a live interaction between AI agents, but to construct a conceptual model of the Linux operating system that is detailed enough to reason about how competing AI agents might behave within it, and to derive from that model a set of strategies and limitations that could inform the design of a future simulation.

In simple terms, the problem is twofold. First, it is unclear which system operations are meaningfully available to an AI agent running locally, and how those operations relate to the agent's ability to monitor, persist, and interfere with a rival. Second, there is currently no established framework for representing the state of a Linux system in a way that captures both the physical layout of data and the strategic position of agents operating within it. This paper attempts to address both by building a layered model, the Operating System Model (OSM) and the Agent Expanded System Model (AESM), that can serve as the foundation for such a simulation.

\subsubsection{Research Question}
\label{ResearchQuestion}

The research question is then as follows:

\begin{center}
    \textit{How can a conceptual model of the Linux operating system be constructed to represent and analyse the interactions between competing AI agents through filesystem operations and process management?}
\end{center}

\begin{itemize}
    \item \textbf{1.} What Linux filesystem and process management commands are available to an AI agent operating locally, and how do they relate to the agent's ability to monitor and interact with another agent on the same system?
    \item \textbf{2.} How can filesystem state, including inode metadata, directory structure, and process information, serve as both a detection surface and an attack vector in AI-to-AI interactions?
    \item \textbf{3.} What theoretical strategies can be derived from the model for agents attempting to achieve operational supremacy over a rival, and what trade-offs do those strategies involve?
    \item \textbf{4.} What are the implications of a filesystem-centric model of AI agent interaction for the design of future autonomous cybersecurity simulations?
\end{itemize}

To address these questions, we will first compile a list of Linux terminal commands that provide the operational foundation for the model. We will then describe the relevant filesystem theory before constructing the OSM and AESM, and finally derive agent strategies and discuss the model's limitations.

\subsection{Scope and Limitations}
\label{ScopeAndLimitations}

This research is intentionally scoped to focus on the Linux operating system and its associated filesystem operations. All commands, system calls, and utilities discussed in this paper relate specifically to the Linux environment and its filesystem-centric architecture. Importantly, this paper produces a preliminary conceptual model rather than an implemented simulation; the OSM and AESM are analytical frameworks intended to structure thinking about agent interactions, not software systems that have been executed or tested.

This limitation has several implications. First, the commands and techniques discussed are primarily applicable to Linux-based systems and may not directly translate to other operating systems such as Windows or macOS, which have fundamentally different filesystem structures, permission models, and command-line interfaces. Second, while Linux systems support various network operations and inter-process communication mechanisms, this study focuses on filesystem-based interactions, as these represent the most fundamental and universally available means of system manipulation and inter-agent communication. Third, the strategies derived from the model are theoretical; their practical viability would need to be validated through an actual simulation before stronger claims could be made.

The rationale for this scope is threefold: (1) Linux provides a well-documented, standardised command-line environment that is widely used in server and embedded systems where autonomous AI agents are most likely to operate; (2) the filesystem serves as a common substrate through which processes can interact, persist data, and observe system state, making it a natural focus for studying AI-to-AI interactions; and (3) limiting the scope to a single operating system allows for deeper analysis of specific mechanisms rather than superficial coverage of multiple platforms.

Consequently, this paper does not address platform-specific considerations for Windows or macOS, nor does it cover network-based attack and defence vectors, which, while important in real-world scenarios, fall outside the scope of this filesystem-focused investigation. 